\section{问题三模型的建立与求解}
\subsection{Cache 查询语义}
问题三沿用问题二的 Task、容量和跨核同步规则，只在所有核心之间增加一个只读 FIFO Cache。对每个符合条件的 COPY\_IN，在发射时按逻辑 tensor id 查询 Cache。命中访问 Cache 带宽，不改变 FIFO 顺序；未命中访问 DDR，完成后按 FIFO 规则插入；容量不足时淘汰最早进入的条目。单个张量大于 Cache 容量时不缓存。Cache 和 DDR 使用独立带宽池，不能将 Cache 命中直接解释为 DDR 搬运或总时间的等比例下降。

设 Cache 服务字节和总查询字节分别为 $Q_L$ 与 $Q_{\rm query}$，字节命中率定义为
\begin{equation}
\rho=\frac{Q_L}{Q_{\rm query}},qquad 0\le\rho\le1.
\end{equation}
完成时间由 DDR 服务、Cache 服务、Pipe 计算、核内容量和跨核同步共同决定。因而问题三的候选目标仍采用官方 Makespan 主排序，并把 $Q^{\rm add}$、$\rho$ 和两类服务池负载作为解释量。

\subsection{快速评分与 FIFO 回放}
轻量层先根据同一张图、同一核数下的 B 方案，计算共享输入在不同消费者分布下的查询次数和可能的驻留关系。随后对候选查询事件进行一次 FIFO 回放，记录查询时间、逻辑键、命中/未命中、Cache 插入和淘汰。轻量评分不替代官方事件模拟；若代理与官方结果方向不一致，保留该候选及其反例，用于修正解释边界。

\subsection{三类候选}
\textbf{（1）聚合候选。}对共享输入按大小与不同消费块数排序，将消费块尽量移动到同一目标核，并在依赖允许的最早和最晚位置尝试插入。该候选降低重复查询，但可能增加同核计算拥挤和私有容量压力。

\textbf{（2）分散候选。}将共享输入的消费块按 Pipe 工作量和依赖关系分配到多个核心，以增加计算并行；同时检查前一次查询是否已完成并写入 Cache。该候选只有在查询事件和 DDR/Cache 服务池的综合评分不恶化时才进入展开。

\textbf{（3）时序候选。}固定分核关系，只调整同核顺序和受依赖允许的查询位置，使重复访问尽可能错开。它独立占用一个局部候选名额，不能因为轻量结构字节不变而被视为没有作用。

三条路径均从同一 B 初始化方案独立开始，每条路径设置固定尝试上限；相同规范化计划只展开一次并记录多个来源。局部、聚合和分散候选按统一计划键进入官方比较，最终选择规则为
\begin{equation}
X_C^*=\arg\min_X\bigl(F_C(X),Q_C^{\rm add}(X),\operatorname{key}(X)\bigr).
\end{equation}

\subsection{问题三的机制对照}
实验完成后，正文将按三个互不混淆的对照回答题目要求：
\begin{enumerate}
\item 固定 B 方案，在无 L2 与有只读 Cache 下比较，衡量硬件资源本身的作用；
\item 在有 Cache 的配置下比较 B 初始化方案与 C 调度适配方案，衡量调度改变的作用；
\item 对不同 Cache 容量和带宽设置做敏感性分析，观察聚合、分散与时序候选的适用条件。
\end{enumerate}
这三项对照的分子、分母、样本和核数必须分别记录，不能把不同方案的最好值拼成一个提升百分比。对应数值、图表和逐用例表在实验完成前均保持 `\pending`。
